\begin{proposition}
\label{proposition-fg-tensor}
Let $M$ be an $R$-module.  The following are equivalent:
\begin{enumerate}
\item $M$ is finitely generated.
\item For every family $(Q_{\alpha})_{\alpha \in A}$ of $R$-modules, the
canonical map $M \otimes_R \left( \prod_{\alpha} Q_{\alpha} \right)
\to \prod_{\alpha} (M \otimes_R Q_{\alpha})$ is surjective.
\item For every $R$-module $Q$ and every set $A$, the canonical map $M
\otimes_R Q^{A} \to (M \otimes_R Q)^{A}$ is surjective.
\item For every set $A$, the canonical map $M \otimes_R R^{A} \to
M^{A}$ is surjective.
\end{enumerate}
\end{proposition}

\begin{proof}
First we prove (1) implies (2).  Choose a surjection $R^n \to M$ and
consider the commutative diagram
$$
\xymatrix{
R^n \otimes_R (\prod_{\alpha} Q_{\alpha})  \ar[r]^{\cong} \ar[d] &
\prod_{\alpha} (R^n \otimes_R Q_{\alpha}) \ar[d] \\
M \otimes_R (\prod_{\alpha} Q_{\alpha})  \ar[r] & \prod_{\alpha} ( M
\otimes_R Q_{\alpha}).
}
$$
The top arrow is an isomorphism and the vertical arrows are surjections.  We
conclude that the bottom arrow is a surjection.

\medskip\noindent
Obviously (2) implies (3) implies (4), so it remains to prove (4) implies (1).
In fact for (1) to hold it suffices that the element $d = (x)_{x \in M}$ of
$M^M$ is in the image of the map $f: M \otimes_R R^{M} \to M^M$.  In
this case $d = \sum_{i = 1}^{n} f(x_i \otimes a_i)$ for some $x_i \in M$ and
$a_i \in R^M$.  If for $x \in M$ we write $p_x: M^M \to M$ for the
projection onto the $x$-th factor, then
$$
x = p_x(d) = \sum\nolimits_{i = 1}^{n} p_x(f(x_i \otimes a_i)) =
\sum\nolimits_{i=1}^{n} p_x(a_i) x_i.
$$
Thus $x_1, \ldots, x_n$ generate $M$.
\end{proof}
